% Written By Enoch Yu
% Downloaded from archive.enochyu.com

\documentclass{article}

\usepackage{graphicx}
\usepackage{amsmath,amsthm,amssymb}
\usepackage[font=small,labelfont=bf]{caption}
\usepackage{tikz}
\usetikzlibrary{calc}
\usepackage{float}
\usepackage[margin=1in]{geometry}
\usepackage{gensymb}
\usepackage{fancyhdr}
\pagestyle{fancy}
\fancyhead[R]{Enoch Yu}
\pagenumbering{gobble}
\usepackage{enumitem}
\newtheorem{theorem}{Theorem}[section]
\newtheorem*{theorem*}{Theorem}
\newtheorem{lemma}[theorem]{Lemma}
\newtheorem{sublemma}{Lemma}[section]
\newtheorem{proposition}{Proposition}
\newtheorem{corollary}{Corollary}[theorem]
\newenvironment{solution}{\begin{trivlist}\item[]{\bf Solution}}{\qed \end{trivlist}}

\pdfinfo{
  /Title (Lucas's Theorem)
}

\title{Lucas's Theorem}
\author{Enoch Yu}
\date{April 2025}

\begin{document}

\section*{Lucas's Theorem}
\begin{theorem*}
    For arbitrary prime $p$ and natural number $n \ge i$, if $n_{(10)}=\overline{n_mn_{m-1}\dots n_1n_0}_{(p)}$ and $i_{(10)}=\overline{i_mi_{m-1}\dots i_1i_0}_{(p)}$ with possible leading zeros, then 
    \[
    _nC_i \equiv \prod_{j=0}^{m} {}_{n_j}C_{i_j} \pmod{p}
    \]
\end{theorem*}

% Beginning of the proof for the theorem
\begin{proof}

%Lemma 0.0.1
\begin{lemma}
        $(a+b)^p \equiv a^p+b^p \pmod{p} \text{ }\forall$ prime $p$ and $a,b\in\mathbb{Z}$.
\end{lemma}
\begin{proof}
    In modular arithmetic, the Freshman's Dream, which states that $(a+b)^2=a^2+b^2$ is not necessarily false. To generalize, the statement $(a+b)^n=a^n+b^n$ may be true in clock arithmetic.
    \\\\
    Using binomial expansion, it is evident that
    \[
    (a+b)^p={}_{p}C_{0}a^{p}b^{0}+{}_{p}C_{1}a^{p-1}b^{1}+\dots+{}_{p}C_{p-1}a^{1}b^{p-1}+{}_{p}C_{p}a^{0}b^{p}
    \]

% Sublemma 0.0.1.1
\begin{sublemma}
    $p|{}_pC_n \text{ }\forall\text{ }0<n<p$.
\end{sublemma}
\begin{proof}
    Let a natural number $N={}_pC_n$. $N$ is also equal to $\frac{p!}{(p-n)!n!}$. However, because $(p-n)!$ and $n!$ are both less than prime $p$, the factor $p$ remains in the denominator when simplified. Thereby, $p$ is a factor in $N$, or ${}_pC_n$.
\end{proof}

\begin{align*}
    \therefore (a+b)^p &\equiv {}_{p}C_{0}a^{p}b^{0}+{}_{p}C_{1}a^{p-1}b^{1}+\dots+{}_{p}C_{p-1}a^{1}b^{p-1}+{}_{p}C_{p}a^{0}b^{p} \pmod{p} \\
    &\equiv {}_{p}C_{0}a^{p}b^{0}+{}_{p}C_{p}a^{0}b^{p} \pmod{p} \\
     &\equiv a^p+b^p \pmod{p}
\end{align*}
\end{proof}

% Corollary 0.0.2.1
\begin{corollary}
    $(1+x)^{p^r} \equiv 1+x^{p^r} \pmod{p} \text{ }\forall$ prime $p$ and $x,r\in\mathbb{N}$.
\end{corollary}
\begin{proof}
    Let $a=1$ and $b=x$.
    \begin{align*}
        (1+x)^p &\equiv 1^p+x^p \pmod{p} \\
        &\equiv 1+x^p \pmod{p}
    \end{align*}
    By exploiting the properties of expansions of the terms, the following conclusion may be driven. Except the first and the last terms, all the terms in between will include ${}_pC_n$ for $0<n<p$.
    \begin{align*}
        (1+x)^{p^{k+1}} &\equiv \left((1+x)^{p}\right)^{p^k} \pmod{p} \\
        &\equiv 1+x^{p^{k+1}} \pmod{p}
    \end{align*}
    Let $r=k+1$ for clarity.
    \[
    \therefore (1+x)^{p^r} \equiv 1+x^{p^r} \pmod{p}
    \]
\end{proof}

\noindent
Let $n=n_mp^m+n'$ and $i=i_mp^m+i'$ by manipulating the base system for $n$ and $i$.
\begin{align*}
    (1+x)^{p^m} &\equiv 1+x^{p^m} \pmod{p} \text{ (By Corollary 0.0.1.1)} \\
    \left((1+x)^{p^m}\right)^{n_m} &\equiv(1+x^{p^m})^{n_m} \pmod{p} \text{ (By property of Modular Arithmetic)} \\
    (1+x)^{n'} \cdot \left((1+x)^{p^m}\right)^{n_m} &\equiv (1+x)^{n'} \cdot (1+x^{p^m})^{n_m} \pmod{p} \\
    (1+x)^n &\equiv (1+x)^{n'} \cdot (1+x^{p^m})^{n_m} \pmod{p}
\end{align*}
Using binomial expansion, it is evident that the following equations are true.
\begin{align*}
    (1+x^{p^m})^{n_m} &= \sum_{j=0}^{n_m} \binom{n_m}{j} 1^{n_m-j}(x^{p^m})^j \\
    &= \sum_{j=0}^{n_m} \binom{n_m}{j} x^{jp^m}\\
    (1+x)^{n'} &= \sum_{k=0}^{n'} \binom{n'}{k} 1^{n'-k}x^k \\
    &= \sum_{k=0}^{n'} \binom{n'}{k} x^k
\end{align*}
Using property of Modular Arithmetic, 
\[(1+x)^n \equiv \left(\sum_{k=0}^{n'} \binom{n'}{k} x^{k}\right) \cdot \left(\sum_{j=0}^{n_m} \binom{n_m}{j} x^{jp^m}\right) \pmod{p}\]
Each term is in the form of $\binom{n'}{k}x^{k}\binom{n_m}{j} x^{jp^m}$, or $\binom{n'}{k}\binom{n_m}{j}x^{k+jp^m}$.
\\ The following expression provides the sum of the coefficient of $x^i$.
\[
\sum_{k+jp^m=i'+i_mp^m}^{}\binom{n'}{k}\binom{n_m}{j}
\]
The reason is because $\binom{n'}{k}\binom{n_m}{j}x^{k+jp^m}$ provides every possible terms with $x^{k+jp^m}$. Thus, replacing $x^{k+jp^m}=x^i$ leads to selected term(s) with $x^i$. The term(s) simply does not exist if $x^{k+jp^m}=x^i$ cannot be true. Recall $i=i'+i_mp^m$. Thereby, the expression above provides the coefficient of $x^i$.
\\\\
The pair(s) of $(k,j)$ that satisfy the equation $k+jp^m=i'+i_mp^m$ must be found.
\[
k-i'+(j-i_m)p^m=0
\]
Due to property of combinations, it is evident that $k \le n'$ and $j \le n_m$. For $k \le n'$, recall that $n'=n_0p^0+n_1p^1+\dots+n_{m-1}p^{m-1}$. The maximum of $n'$ is when $n_0,n_1,\dots n_{m-1}$ are maximum, or $p-1$. In another words,
\begin{align*}
\max(n')&=(p-1)p^0+(p-1)p^1+\dots+(p-1)p_{m-1} \\
&=(p-1)(p^0+p^1+\dots p^{m-1}) \\
&=(p-1)[p^0\cdot\frac{1-p^m}{1-p}] \\
&=p^m-1.
\end{align*}
Because $\max(n')<p^m$, $n'<p^m$.
\\\\
For $j \le n_m$, Using the properties of base number representation, it is evident that $j<p$ because $n_m<p$.
\end{proof}


\section*{Example}
\textbf{Problem} Math Virus School has a total of 4 classes, and each class contains a certain number of viruses: 1246, 896, 4050, and 7547 viruses, respectively. Every virus from all the classes leaves its class and selects one hat from 58 distinct hats, each differing in shape and color. After making a selection, the viruses return to their respective classes. Let the remainder when the total number of possible selections is divided by 7 be $r_7$, and the remainder when it is divided by 3 be $r_3$. Find the value $r_7 × r_3$.
\\\\
\textbf{Key Word} Lucas's Theorem, Stars and Bars, Modular Arithmetic Properties
\\\\
Let stars represent the number of viruses in each class for each case. If 57 bars are used, the number of cases in which each virus choose one hat in each class could be obtained.
\\\\
\begin{tabular}{c|c|c|c}
    \boxed{1246} & \boxed{896} & \boxed{4050} & \boxed{7547} \\
    \text{1246 Stars, 57 Bars} & \text{896 Stars, 57 Bars} & \text{4050 Stars, 57 Bars} & \text{7547 Stars, 57 Bars} \\
    $\frac{(1246+57)!}{1246!57!}$ & $\frac{(896+57)!}{896!57!}$ & $\frac{(4050+57)!}{4050!57!}$ & $\frac{(7547+57)!}{7547!57!}$ \\
    $={}_{1303}C_{57}$ & $={}_{953}C_{57}$ & $={}_{4107}C_{57}$ & $={}_{7604}C_{57}$ \\
\end{tabular}
\\\\
The total probability is ${}_{1303}C_{57} \cdot {}_{953}C_{57} \cdot {}_{4107}C_{57} \cdot {}_{7604}C_{57}$
\\\\
Finding remainder when a number involving multiplication of combinations is divided by an integer reminds Lucas's Theorem.
\\\\
\textbf{Applying Lucas's Theorem for $r_7$}
\begin{align*}
57_{(10)}&=111_{(7)} \\
1303_{(10)}&=3542_{(7)} \\
953_{(10)}&=2531_{(7)} \\
4107_{(10)}&=14655_{(7)} \\
7604_{(10)}&=31111_{(7)} \\
\end{align*}
Lucas's Theorem is applied to find $r_7$.
\begin{align*}
    {}_{1303}C_{57}
    &\equiv {}_{3}C_{0}\cdot{}_{5}C_{1}\cdot{}_{4}C_{1}\cdot{}_{2}C_{1} \pmod{7} \\
    &\equiv 1\cdot5\cdot4\cdot2 \pmod{7} \\
    &\equiv 40 \pmod{7} \\
    &\equiv 5 \pmod{7}
    \\\\
    {}_{953}C_{57}
    &\equiv {}_{2}C_{0}\cdot{}_{5}C_{1}\cdot{}_{3}C_{1}\cdot{}_{1}C_{1} \pmod{7} \\
    &\equiv 1\cdot5\cdot3\cdot1 \pmod{7} \\
    &\equiv 15 \pmod{7} \\
    &\equiv 1 \pmod{7}
\end{align*}
\begin{align*}
    {}_{4107}C_{57}
    &\equiv {}_{1}C_{0}\cdot{}_{4}C_{0}\cdot{}_{6}C_{1}\cdot{}_{5}C_{1}\cdot{}_{5}C_{1} \pmod{7} \\
    &\equiv 1\cdot1\cdot6\cdot5\cdot5 \pmod{7} \\
    &\equiv 150 \pmod{7} \\
    &\equiv 3 \pmod{7}
    \\\\
    {}_{7604}C_{57}
    &\equiv {}_{3}C_{0}\cdot{}_{1}C_{0}\cdot{}_{1}C_{1}\cdot{}_{1}C_{1}\cdot{}_{1}C_{1} \pmod{7} \\
    &\equiv 1\cdot1\cdot1\cdot1\cdot1 \pmod{7} \\
    &\equiv 1 \pmod{7}  
\end{align*}
\begin{align*}
    {}_{1303}C_{57}\cdot{}_{953}C_{57}\cdot{}_{4107}C_{57}\cdot{}_{7604}C_{57} &\equiv 5\cdot1\cdot3\cdot1 \pmod{7} \\
    &\equiv 15 \pmod{7} \\
    &\equiv 1 \pmod{7}
\end{align*}
$\therefore r_7=1$
\\\\
\textbf{Applying Lucas's Theorem for $r_3$}
\begin{align*}
57_{(10)}&=2010_{(3)} \\
1303_{(10)}&=1201021_{(3)} \\
953_{(10)}&=1022022_{(3)} \\
4107_{(10)}&=111212100_{(3)} \\
7604_{(10)}&=101102122_{(3)} \\
\end{align*}
\begin{align*}
    {}_{1303}C_{57}
    &\equiv {}_{1}C_{0}\cdot{}_{2}C_{0}\cdot{}_{0}C_{0}\cdot{}_{1}C_{2}\cdot{}_{0}C_{0}\cdot{}_{2}C_{1}\cdot{}_{1}C_{0} \pmod{3} \\
    &\equiv 1\cdot1\cdot1\cdot0\cdot1\cdot2\cdot1 \pmod{3} \\
    &\equiv 0 \pmod{3}
\end{align*}
$r_3=0 \text{ }(\because {}_{1303}C_{57} \equiv0 \pmod{3})$ \\
\\\\
$r_7 \cdot r_3=1\cdot0=0$
\\\\
\\\\
\textbf{Conclusion} $r_7 \cdot r_3=\boxed{0}$.

\end{document}

